\section{Obiekty rzutowe, injektywne i rezolucje}\label{cat:sec:projective-injective}
Pracujemy w lokalnie małej kategorii abelowej $\mathcal A$. Wykorzystujemy teorię funktorów dokładnych z rozdziału~\ref{cat:sec:abelian-categories}. W przypadku modułów odpowiednie pojęcie rzutowości oraz jego charakterystykę przez składniki sum prostych modułów wolnych znamy z \atlasref{AA-16::aa:tw:modules-projective-summand}{AA-16: charakterystyka modułów rzutowych}.

\subsection{Rzutowość i injektywność}
\begin{definicja}\label{cat:def:projective-injective}
Obiekt $P$ jest rzutowy, jeżeli dla każdego epimorfizmu $q:E\to N$ i każdego $f:P\to N$ istnieje $\widetilde f:P\to E$ z $q\widetilde f=f$.
Obiekt $I$ jest injektywny, jeżeli dla każdego monomorfizmu $j:M\to E$ i każdego $g:M\to I$ istnieje $\widetilde g:E\to I$ z $\widetilde g j=g$. Nie wymagamy jednoznaczności tych morfizmów.
\[
\begin{array}{ccc}
P&&\\
\downarrow\widetilde f&\searrow f&\\
E&\xrightarrow{q}&N
\end{array}
\qquad
\begin{CD}
M @>{j}>> E\\
@V{g}VV @VV{\widetilde g}V\\
I @>{\id_I}>> I
\end{CD}
\]
\end{definicja}
\begin{twierdzenie}[Charakteryzacje rzutowości i injektywności]\label{cat:thm:projective-characterization}
Dla obiektu $P$ równoważne są następujące warunki:
\begin{enumerate}
\item $P$ jest rzutowy;
\item funktor $\Hom(P,-):\mathcal A\to\Ab$ jest dokładny;
\item każdy epimorfizm $E\to P$ ma prawą odwrotność;
\item każdy krótki ciąg dokładny kończący się obiektem $P$ jest rozkładalny.
\end{enumerate}
Dualnie dla obiektu $I$ równoważne są injektywność, dokładność $\Hom(-,I):\mathcal A^{\op}\to\Ab$, istnienie lewej odwrotności każdego monomorfizmu $I\to E$ oraz rozkładalność każdego krótkiego ciągu dokładnego rozpoczynającego się obiektem $I$.
\end{twierdzenie}
\begin{proof}
Funktor $\Hom(P,-)$ jest lewostronnie dokładny na mocy przykładu~\ref{cat:ex:hom-left-exact}. Dla krótkiego ciągu dokładnego $0\to A\to E\xrightarrow{q}N\to0$ brakuje zatem jedynie surjektywności $\Hom(P,q)$. Jest ona równoważna istnieniu podniesienia każdego $f:P\to N$ wzdłuż $q$. Każdy epi $q$ występuje w takim ciągu, ponieważ jest kojądrem swego jądra. To daje równoważność (1) i (2).
Warunek (1) zastosowany do $f=\id_P$ daje (3). Jeśli zachodzi (3), to dla dowolnego epi $q:E\to N$ i $f:P\to N$ tworzymy iloczyn włóknisty $D=E\times_N P$ z projekcjami $a:D\to E$, $b:D\to P$. Z twierdzenia~\ref{cat:thm:abelian-base-change} morfizm $b$ jest epi. Istnieje więc $s:P\to D$ z $bs=\id_P$. Morfizm $\widetilde f=as$ spełnia $q\widetilde f=qas=fbs=f$. Otrzymujemy (1). Równoważność (3) i (4) wynika z twierdzenia~\ref{cat:thm:abelian-split} i z tego, że każdy epimorfizm $E\to P$ jest kojądrem swego jądra. Twierdzenie dla $I$ otrzymujemy w kategorii przeciwnej, która jest abelowa.
\end{proof}
\begin{stwierdzenie}[Retrakty i sumy proste]\label{cat:prop:projective-retracts}
Retrakt obiektu rzutowego jest rzutowy. Dowolny istniejący koprodukt rodziny obiektów rzutowych jest rzutowy. Dualnie retrakty obiektów injektywnych i dowolne istniejące iloczyny obiektów injektywnych są injektywne.
\end{stwierdzenie}
\begin{proof}
Niech $i:P\to Q$, $r:Q\to P$ spełniają $ri=\id_P$, gdzie $Q$ jest rzutowy. Dla epi $q:E\to N$ i $f:P\to N$ podnosimy $fr:Q\to N$ do $h:Q\to E$. Wtedy $qhi=fri=f$. Jeśli $P=\coprod_{s\in S}P_s$ z włożeniami $i_s$, wybieramy podniesienia $h_s:P_s\to E$ morfizmów $fi_s$. Dla nieskończonego zbioru $S$ korzystamy z aksjomatu wyboru. Własność koproduktu daje $h:P\to E$ z $hi_s=h_s$. Ponieważ $qhi_s=fi_s$ dla każdego $s$, mamy $qh=f$. Wnioski o injektywności są dualne.
\end{proof}
\begin{przyklad}
W $\Ab$ obiekt $\mathbb Z$ jest rzutowy, ponieważ homomorfizm z $\mathbb Z$ wyznacza obraz $1$ i można go podnieść wzdłuż każdej surjekcji. Obiekt $\mathbb Z/2\mathbb Z$ nie jest rzutowy, ponieważ ciąg $0\to\mathbb Z\xrightarrow{2}\mathbb Z\to\mathbb Z/2\mathbb Z\to0$ nie jest rozkładalny. Nie jest też injektywny: homomorfizm $2\mathbb Z\to\mathbb Z/2\mathbb Z$ wysyłający $2$ na klasę $1$ nie przedłuża się do $\mathbb Z$, gdyż każde przedłużenie $g$ musiałoby spełniać $g(2)=2g(1)=0$. Podobnie $\mathbb Z$ nie jest injektywna: homomorfizm $2\mathbb Z\to\mathbb Z$ wysyłający $2$ na $1$ nie może się przedłużyć, gdyż równanie $2g(1)=1$ nie ma rozwiązania w $\mathbb Z$.
\end{przyklad}

\subsection{Dostatecznie wiele obiektów i rezolucje}
\begin{definicja}\label{cat:def:enough-projectives-injectives}
Mówimy, że kategoria abelowa ma dostatecznie wiele obiektów rzutowych, gdy dla każdego $X$ istnieje epi $P\to X$ z $P$ rzutowym. Ma dostatecznie wiele obiektów injektywnych, gdy dla każdego $X$ istnieje mono $X\to I$ z $I$ injektywnym.
\end{definicja}
\begin{przyklad}[Moduły]\label{cat:ex:enough-projectives-mod}
Kategoria lewych modułów nad dowolnym pierścieniem $R$ z jedynką ma dostatecznie wiele obiektów rzutowych: moduł wolny $R^{(M)}$ na zbiorze elementów $M$ jest rzutowy, a homomorfizm wysyłający generator indeksowany przez $m$ na $m$ jest surjektywny. Rzutowość modułu wolnego oraz własność uniwersalną tej konstrukcji udowodniono w \atlasref{AA-16::aa:tw:modules-projective-summand}{AA-16: moduły rzutowe} oraz \atlasref{AA-16::aa:tw:modules-free-universal}{AA-16: moduły wolne}. W kategorii abelowej mającej dostatecznie wiele obiektów rzutowych każdy obiekt rzutowy jest retraktem każdego obiektu rzutowego, z którego istnieje na niego epimorfizm, na mocy twierdzenia~\ref{cat:thm:projective-characterization}.
\end{przyklad}
\begin{definicja}[Rezolucje]\label{cat:def:resolutions}
Rezolucją rzutową obiektu $X$ nazywamy dokładny ciąg
\[
 \cdots\xrightarrow{d_3}P_2\xrightarrow{d_2}P_1
 \xrightarrow{d_1}P_0\xrightarrow{\epsilon}X\to0,
\]
w którym wszystkie $P_n$ są rzutowe. Rezolucja injektywna ma postać dokładnego ciągu
\[
 0\to X\xrightarrow{\eta}I^0\xrightarrow{d^0}I^1
 \xrightarrow{d^1}I^2\to\cdots,
\]
w którym wszystkie $I^n$ są injektywne. Dokładność wymusza zerowanie złożeń dwóch kolejnych strzałek.
\end{definicja}
\begin{twierdzenie}[Istnienie rezolucji]\label{cat:thm:resolution-existence}
Jeżeli kategoria abelowa ma dostatecznie wiele obiektów rzutowych, to każdy obiekt ma rezolucję rzutową. Jeżeli ma dostatecznie wiele obiektów injektywnych, to każdy obiekt ma rezolucję injektywną. Przy nieskończonej konstrukcji przyjmujemy możliwość kolejnych wyborów obiektów i strzałek (na przykład pracujemy względem ustalonego uniwersum i stosujemy aksjomat wyboru zależnego).
\end{twierdzenie}
\begin{proof}
Przyjmijmy $K_0=X$. Wybieramy epi $e_0:P_0\to K_0$ z obiektem rzutowym $P_0$ i kładziemy $\epsilon=e_0$. Indukcyjnie, gdy mamy $e_n:P_n\to K_n$, wybieramy jądro $k_{n+1}:K_{n+1}\to P_n$ oraz epi $e_{n+1}:P_{n+1}\to K_{n+1}$ z obiektem rzutowym $P_{n+1}$. Definiujemy $d_{n+1}=k_{n+1}e_{n+1}$. Ponieważ $e_{n+1}$ jest epi, a $k_{n+1}$ mono, obraz $d_{n+1}$ jest reprezentowany przez $k_{n+1}$: ich kojądra mają tę samą własność uniwersalną, gdyż $hd_{n+1}=0$ wtedy i tylko wtedy, gdy $hk_{n+1}=0$.
Dla $n\ge1$ zachodzi $\ker d_n=\ker e_n$, bo $k_n$ jest mono; dla $n=0$ mamy $\ker\epsilon=\ker e_0$. W każdym przypadku $k_{n+1}$ jest więc jądrem strzałki wychodzącej z $P_n$. Tym samym otrzymujemy dokładność w każdym $P_n$, a dokładność w $X$ wynika z epimorficzności $\epsilon$. Rezolucję injektywną konstruujemy stosując ten dowód w $\mathcal A^{\op}$.
\end{proof}
\begin{przyklad}[Krótka rezolucja rzutowa]\label{cat:ex:cyclic-projective-resolution}
Dla $m\ge2$ ciąg $0\to\mathbb Z\xrightarrow{m}\mathbb Z\to\mathbb Z/m\mathbb Z\to0$ jest rezolucją rzutową grupy $\mathbb Z/m\mathbb Z$ (w stopniach $n\ge2$ przyjmujemy $P_n=0$). Nie jest to rozkładalny ciąg dokładny: gdyby projekcja miała prawą odwrotność, obraz klasy $1$ byłby elementem $\mathbb Z$ zabijanym przez $m$, czyli zerem, co jest niemożliwe.
\end{przyklad}

\subsection{Porównywanie rezolucji i homotopie}
Dla kompleksu łańcuchowego $P_\bullet$ zapisujemy różniczki jako $d_n^P:P_n\to P_{n-1}$, $n\ge1$, z $d_n^Pd_{n+1}^P=0$. Morfizm kompleksów $f:P_\bullet\to Q_\bullet$ jest rodziną $f_n:P_n\to Q_n$ spełniającą $d_n^Q f_n=f_{n-1}d_n^P$. W rezolucjach obiektów $X,Y$ morfizm ten podnosi $u:X\to Y$, gdy $\epsilon_Q f_0=u\epsilon_P$. W przypadku modułów definicje omówiono w \atlasref{TOP-03::def:kompleks-lancuchowy}{TOP-03: kompleksy łańcuchowe} i \atlasref{TOP-03::def:mapa-lancuchowa}{TOP-03: morfizmy łańcuchowe}.
\begin{definicja}[Homotopia łańcuchowa]\label{cat:def:resolution-homotopy}
Morfizmy kompleksów $f,g:P_\bullet\to Q_\bullet$ są homotopijne, jeśli istnieją $h_n:P_n\to Q_{n+1}$ takie, że
\begin{equation}\label{cat:eq:resolution-homotopy}
 f_0-g_0=d_1^Qh_0,\qquad
 f_n-g_n=d_{n+1}^Qh_n+h_{n-1}d_n^P\quad(n\ge1).
\end{equation}
Homotopia jest relacją równoważności: homotopia zerowa, zmiana znaku i dodawanie homotopii dają zwrotność, symetrię i przechodniość. Por. \atlasref{TOP-03::def:homotopia-lancuchowa}{TOP-03: homotopia łańcuchowa}.
\end{definicja}
\begin{twierdzenie}[Twierdzenie porównawcze]\label{cat:thm:resolution-comparison}
Niech $P_\bullet\to X$ będzie rezolucją rzutową, a $Q_\bullet\to Y$ dokładnym ciągiem augmentowanym postaci z definicji~\ref{cat:def:resolutions}, bez założenia rzutowości $Q_n$. Każdy $u:X\to Y$ podnosi się do morfizmu kompleksów $f:P_\bullet\to Q_\bullet$. Dowolne dwa podniesienia są homotopijne.
\end{twierdzenie}
\begin{proof}
W stopniu zerowym podnosimy $u\epsilon_P$ wzdłuż epimorfizmu $\epsilon_Q$ dzięki rzutowości $P_0$. Otrzymujemy $f_0$ z $\epsilon_Q f_0=u\epsilon_P$. Załóżmy, że skonstruowano $f_0,\ldots,f_{n-1}$ zgodne z różniczkami i augmentacją. Dla $n=1$ mamy $\epsilon_Q f_0d_1^P=u\epsilon_Pd_1^P=0$, a dla $n\ge2$
\[
 d_{n-1}^Q f_{n-1}d_n^P=f_{n-2}d_{n-1}^Pd_n^P=0.
\]
Niech $k:K\to Q_{n-1}$ będzie jądrem $\epsilon_Q$ dla $n=1$ lub $d_{n-1}^Q$ dla $n\ge2$. Istnieje jedyny $a:P_n\to K$ z $ka=f_{n-1}d_n^P$. Z dokładności $Q$ różniczka $d_n^Q$ rozkłada się jako $ke$ dla epi $e:Q_n\to K$. Rzutowość $P_n$ daje $f_n:P_n\to Q_n$ z $ef_n=a$, a stąd $d_n^Qf_n=ka=f_{n-1}d_n^P$. Indukcja daje istnienie.

Niech $f,g$ podnoszą ten sam $u$. Ponieważ $\epsilon_Q(f_0-g_0)=0$, morfizm $f_0-g_0$ rozkłada się przez jądro $\epsilon_Q$. Dokładność $Q$ i rzutowość $P_0$ pozwalają go podnieść do $h_0:P_0\to Q_1$ z $d_1^Qh_0=f_0-g_0$.
Jeśli skonstruowano $h_0,\ldots,h_{n-1}$ spełniające równania w niższych stopniach, kładziemy
\[
 a_n=f_n-g_n-h_{n-1}d_n^P:P_n\to Q_n.
\]
Dla $n=1$ otrzymujemy
$d_1^Qa_1=(f_0-g_0)d_1^P-d_1^Qh_0d_1^P=0$.
Dla $n\ge2$
\[
\begin{aligned}
 d_n^Qa_n
 &=(f_{n-1}-g_{n-1}-d_n^Qh_{n-1})d_n^P\\
 &=h_{n-2}d_{n-1}^Pd_n^P=0.
\end{aligned}
\]
Zatem $a_n$ rozkłada się przez jądro $d_n^Q$. Dokładność $Q$ dostarcza epimorfizmu z $Q_{n+1}$ na to jądro; rzutowość $P_n$ daje $h_n$ z $d_{n+1}^Qh_n=a_n$. Jest to równanie~\eqref{cat:eq:resolution-homotopy}. Indukcja kończy dowód (dla nieskończonych ciągów korzystamy z kolejnych wyborów jak przy istnieniu rezolucji).
\end{proof}
\begin{wniosek}[Niezależność rezolucji]\label{cat:cor:resolution-homotopy-equivalence}
Dwie rezolucje rzutowe tego samego obiektu są równoważne homotopijnie. Dla ustalonego wyboru rezolucji obiektów podniesienie morfizmu do klasy homotopii jest funktorialne.
\end{wniosek}
\begin{proof}
Stosujemy twierdzenie~\ref{cat:thm:resolution-comparison} do $\id_X$ w obu kierunkach, otrzymując $f:P_\bullet\to Q_\bullet$ i $g:Q_\bullet\to P_\bullet$. Złożenia $gf,fg$ podnoszą $\id_X$, podobnie jak identyczności kompleksów, więc $gf\simeq\id_P$ oraz $fg\simeq\id_Q$. Jeżeli $f,g$ podnoszą kolejno $u:X\to Y$, $v:Y\to Z$, to $\epsilon_Rg_0f_0=v\epsilon_Qf_0=vu\epsilon_P$. Stąd złożenie klas podniesień jest klasą podniesienia $vu$, a klasa podniesienia identyczności jest klasą identyczności.
Składanie jest dobrze określone na klasach homotopii: jeśli $f-f'=d^Qh+hd^P$ i $a:Q_\bullet\to R_\bullet$ jest morfizmem kompleksów, to
$af-af'=d^R(ah)+(ah)d^P$, gdzie $(ah)_n=a_{n+1}h_n$, dzięki przemienności $a$ z różniczkami.
Dla $b:S_\bullet\to P_\bullet$ mamy
$fb-f'b=d^Q(hb)+(hb)d^S$, gdzie $(hb)_n=h_nb_n$, dzięki przemienności $b$ z różniczkami.
Zatem zastąpienie obu czynników morfizmami homotopijnymi nie zmienia klasy złożenia.
\end{proof}
\begin{wniosek}[Porównywanie rezolucji injektywnych]\label{cat:cor:injective-comparison}
Niech $X\to I^\bullet$ będzie dokładnym ciągiem koaugmentowanym, a $Y\to J^\bullet$ rezolucją injektywną. Każdy $u:X\to Y$ przedłuża się do morfizmu kompleksów $f:I^\bullet\to J^\bullet$, jednoznacznego z dokładnością do homotopii kołańcuchowej. Dwie rezolucje injektywne tego samego obiektu są równoważne homotopijnie.
\end{wniosek}
\begin{proof}
Stosujemy twierdzenie~\ref{cat:thm:resolution-comparison} i wniosek~\ref{cat:cor:resolution-homotopy-equivalence} w kategorii przeciwnej. Tam $J^\bullet$ staje się rezolucją rzutową, a $u$ ma kierunek od $Y$ do $X$. Po powrocie do $\mathcal A$ otrzymujemy morfizm we wskazanym kierunku. Homotopia ma postać $s^n:I^n\to J^{n-1}$ dla $n\ge1$ i spełnia
\[
 f^n-g^n=d_J^{n-1}s^n+s^{n+1}d_I^n,
\]
gdzie pierwszy składnik dla $n=0$ jest zerem. Jest to odwrócenie złożeń z~\eqref{cat:eq:resolution-homotopy}.
\end{proof}

\subsection{Lemat Schanuela}
\begin{twierdzenie}[Lemat Schanuela]\label{cat:thm:schanuel}
Niech
\[
 0\to K\xrightarrow{i}P\xrightarrow{p}X\to0,\qquad
 0\to L\xrightarrow{j}Q\xrightarrow{q}X\to0
\]
będą krótkimi ciągami dokładnymi, w których $P,Q$ są rzutowe. Wówczas $K\oplus Q\cong L\oplus P$.
\end{twierdzenie}
\begin{proof}
Tworzymy iloczyn włóknisty $D=P\times_X Q$ z projekcjami $a:D\to P$, $b:D\to Q$. Obie projekcje są epi na mocy twierdzenia~\ref{cat:thm:abelian-base-change}. Para $(i,0)$ definiuje $k:K\to D$ z $ak=i$, $bk=0$. Jeśli $z:T\to D$ spełnia $bz=0$, to $paz=qbz=0$, więc istnieje jedyny $y:T\to K$ z $iy=az$. Morfizmy $ky$ i $z$ mają takie same złożenia z $a,b$, a zatem są równe przez własność iloczynu włóknistego. Jednoznaczność $y$ wynika z monomorficzności $i$. To dowodzi, że $k=\ker b$. Analogicznie para $(0,j)$ definiuje jądro $l:L\to D$ morfizmu $a$: dla $az=0$ mamy $qbz=0$, zatem $bz=jy$ dla jedynego $y:T\to L$ i $ly=z$.
Otrzymujemy krótkie ciągi dokładne
\[
 0\to K\xrightarrow{k}D\xrightarrow{b}Q\to0,\qquad
 0\to L\xrightarrow{l}D\xrightarrow{a}P\to0.
\]
Rzutowość $P,Q$ daje przekroje obu epimorfizmów, więc twierdzenie~\ref{cat:thm:abelian-split} daje $D\cong K\oplus Q$ oraz $D\cong L\oplus P$. Składając jeden izomorfizm z odwrotnością drugiego, otrzymujemy tezę. Wybór przekrojów nie musi być jednoznaczny, dlatego uzyskany izomorfizm nie jest na ogół kanoniczny.
\end{proof}
\begin{wniosek}\label{cat:cor:schanuel-projectivity}
Jeżeli $X$ ma krótką rezolucję rzutową $0\to K\to P\to X\to0$ z $K$ rzutowym, to jądro każdego epimorfizmu $Q\to X$ z $Q$ rzutowym jest rzutowe.
\end{wniosek}
\begin{proof}
Dla jądra $L$ tego epimorfizmu lemat Schanuela daje $L\oplus P\cong K\oplus Q$. Prawa strona jest rzutowa jako skończony koprodukt obiektów rzutowych. Obiekt $L$ jest retraktem lewej strony, zatem jest rzutowy na mocy stwierdzenia~\ref{cat:prop:projective-retracts}.
\end{proof}
